\section{Методы}
\label{sec:methods}

\subsection{Участники}

Данные собраны онлайн через русскоязычный DEIC-Inventory. Общая выборка $N = 1426$ после удаления некорректных ответов (менее 80\% заполненности, дублированные идентификаторы, implausible response patterns). Рекрутинг --- self-selected volunteers через Telegram-каналы, VK-сообщества, личные сети исследователя. Инклюзивные критерии: возраст $\geq 18$, русский как язык опросника (non-native speakers не исключались, но не таргетировались). Эксклюзивных критериев не вводилось.

Участие анонимное, без идентификаторов, без оплаты. Добровольное информированное согласие получено на первом экране опросника. Этическое одобрение: \textit{[to fill]}.

Ограничения выборки (адресуются в §\ref{sec:limitations}): self-selection bias в сторону психологически любопытной, рефлексивной аудитории; язык-специфичность; возрастной перекос в сторону 18--37; отсутствие клинических диагностических оценок (все диагнозы --- self-reported).

\subsection{Инструмент}

DEIC-Inventory --- 131-вопросный self-report инструмент, разработанный итеративно через несколько волн качественной pilot-работы и количественной ревизии. Из них 121 --- классические Ликертовские айтемы; 10 --- ситуационно-сценарные (multiple-choice), собираемые для дескриптивного использования и не включаемые в CFA/latent-анализ. Айтемы распределены по 14 scale-меткам: эмпатия (E), психопатия (P), подавление (S), диссоциация (D), аффективная изоляция (AIS), притупление чувств (Flattening), внутренний диссонанс (IDS), маскирование (M), детский негативный опыт (ACE), детское отчуждение (CA), семейный эмоциональный климат (EC), индекс аффективного нарушения (ADI), внимание (A), контрольная шкала лжи (Lie).

Шкала ответа Ликертовских айтемов: 5-балльная (1 = «совершенно не согласен» $\to$ 5 = «полностью согласен»). Часть айтемов reverse-coded для контроля response bias. Девять Ликертовских айтемов (S\_1--S\_4, S\_new\_1, AIS\_3--AIS\_5, AIS\_new\_5) сохранены в текущем анализе для сопоставимости с историческими отчётами волны~1, но помечены для удаления в волне~2 на основании низких коммуналитетов и содержательного дублирования; их исключение не меняет знак и направление ни одного репортируемого эффекта.

Инновационные особенности инструмента: разделение \emph{декларативных} и \emph{поведенческих} формулировок в парных айтемах (пример: «я считаю, что эмпатия важна» --- декларативный; «я часто отменяю встречи, чтобы побыть с плачущим другом» --- поведенческий); контроль структурных артефактов через балансирование направления шкалы ответа и reverse-coding; Lie-шкала, используемая как контрольная ковариата в partial correlation анализах. Полная item-level спецификация и веса доступны в \texttt{questions.json}.

\paragraph{Техническая оговорка по Lie-ковариате.} В wave-1 Lie операционализирована единственным айтемом (LIE\_1: «Я иногда говорю неправду») с reverse-coding. Одно-айтемная формулировка не позволяет вычислить внутреннюю согласованность и представляет собой методологическое ограничение, признаваемое в §\ref{sec:limitations}; Lie используется face-валидно как индикатор самопрезентационного давления, но без ω-оценки. Wave-2 расширяет Lie до $\geq 6$ айтемов с reported reliability (§\ref{subsec:wave2_design}).

\paragraph{Проверка внимания и QC-пайплайн.} В инструмент встроен один attention-check айтем (CONTROL\_1 с директивной инструкцией выбрать «Скорее согласен»). Респонденты с провалом attention-check флагируются и исключаются из latent-анализов; финальный $N = 1426$ получен из сырого $N = 1470$ после применения QC-фильтров (attention-check, заполненность $\geq 80\%$, implausible response patterns, дедупликация по fingerprint, фильтрация $e2eTestRunId$).

\paragraph{Три параллельных слоя оценки.} Экспортируемый датасет содержит \texttt{raw\_*} (сырая сумма до стандартизации), \texttt{score\_*} (теоретически взвешенный композит с multi-factor-весами из \texttt{questions.json}) и \texttt{mapped\_*} (нормализованная шкала 0--100 для клинического отображения). Все анализы настоящей статьи используют слой \texttt{score\_*}; \texttt{raw\_*} и \texttt{mapped\_*} сохраняются как replication-convenience и как вход для внешних инструментов соответственно.

\paragraph{Продуктовые слои, не входящие в формальный анализ.} Инструмент в рамках продуктовой инфраструктуры дополнительно присваивает каждому респонденту (a) архетипный ярлык из 12 категорий (Трикстер, Странник, Аналитик, \dots) и (b) rule-based композитный профиль-тэг (\texttt{psychometrics.json.profiles}: напр., empathicParadox, defensive\_hardness, modestlyBlocked). Ни один из этих слоёв не входит в latent-анализы данной статьи; они сохраняются в экспорте (\texttt{cluster}, \texttt{profiles}, \texttt{ratio}) как задел для отдельной applied-публикации. Вспомогательное поле \texttt{score\_maskarading} --- legacy-дубликат \texttt{mapped\_masking}, остающийся в wave-1 экспорте и запланированный к удалению в wave-2.

\paragraph{Оговорка об item-count heterogeneity.} Семь записей имеют нетипичное количество ответов (121, 124, 127, 130 вместо канонических 122). Это \emph{валидные} ответы ранних pilot-версий инструмента, собранные до того, как item-набор стабилизировался в финальной конфигурации из 122 активных айтемов. Эти респонденты \emph{включены} в аналитическую выборку наравне с остальными; на item-уровне пункты, отсутствовавшие в их версии инструмента, обрабатываются стандартным missing-data handling (listwise для сырых композитов с full-coverage-критерием, FIML-eq для CFA-моделей).

\subsection{Операционализация шести факторов}
\label{subsec:six_factor_ops}

Текущая статья работает с шестью содержательными факторами из четырнадцати шкал инструмента. Ниже --- как каждый из них сконструирован для анализа.

\subsubsection{Факторы через 6-factor oblique CFA}

Первичное операционализирование --- через латентные факторы из шестифакторной эксплоративно-конфирматорной последовательности:
\begin{enumerate}
    \item EFA с parallel analysis на split-half A ($N = 713$) определила $k = 6$ как оптимальное число факторов.
    \item 6-factor oblique CFA (maximum likelihood, correlated factors) на split-half B ($N = 713$) дала $\chi^2(120) = 516.04$, $p < .001$, CFI~=~0.950, TLI~=~0.936, RMSEA~=~0.068 (robust bootstrap approximation), $\chi^2/df = 4.30$.
    \item Latent factor scores экстрагированы через empirical Bayes regression method.
\end{enumerate}

Факторы: \textbf{ID} (Internal Dissonance), \textbf{CA} (Childhood Alienation), \textbf{P} (Psychopathy), \textbf{M} (Masking), \textbf{\EAF{}} (Affective Empathy), \textbf{\Ecog{}} (Cognitive Empathy).

Reliability (McDonald's $\omega$ на latent factor scores): ID~=~0.91, CA~=~0.89, P~=~0.93, \EAF{}~=~0.90, \Ecog{}~=~0.89, \textbf{M~=~0.546}. M-надёжность подчёркнуто низка; M-специфические выводы в статье не делаются. M сохранён в модели как координата для будущей переработки.

\subsubsection{Attention (A): формативный композит}

A не является одним из шести факторов из 14-шкального инструмента. A сконструирован как формативный композит по теоретической схеме: 31 айтем из 9 шкал (E, IDS, D, ADI, AIS, S, EC, CA, P) с a-priori коэффициентами от $+0.4$ до $-0.4$, отражающими теоретическое положение каждого айтема относительно witnessing-функции. Композит:
\[
A_i = \sum_{j=1}^{31} w_j \cdot x_{ij}
\]
где $w_j$ --- теоретический коэффициент айтема, $x_{ij}$ --- $z$-стандартизованный ответ участника $i$ на айтем $j$. Айтемы с отрицательными коэффициентами инвертируют свой вклад, так что высокий балл A означает повышенное witnessing-присутствие.

Коэффициенты установлены на основе итеративной ревизии с участием автора; подробности --- \texttt{validate\_attention.py}. В §\ref{sec:results}.7 мы эмпирически проверяем, что этот формативный композит не может быть редуцирован к единому рефлексивному латентному фактору.

\subsubsection{Альтернативные конструкции (контроль)}

Параллельно с латентными факторами использовались \emph{weighted composites} (классическая реализация через единичные веса по primary\_scale). Эти композиты инфлированы относительно латентных (средняя инфляция $\sim 0.18$, максимум $0.68$ с переменой знака P$\leftrightarrow$E); они приводятся для сравнения, но все содержательные выводы статьи опираются на латентные factor scores.

\subsection{Аналитическая стратегия: прочтение для не-психометриков}
\label{subsec:analytic_primer}

Эта секция написана для читателя, который пришёл в статью из клиники, теории, философии или созерцательной традиции и не обязан свободно владеть жаргоном структурного моделирования. Специалист в психометрии может перейти сразу к §\ref{subsec:stat_approach}. Здесь --- короткое, содержательно-точное объяснение того, что именно делает каждый из анализов, на которые опирается статья, и чего в них \emph{нельзя} увидеть.

\subsubsection{Латентный фактор: почему не просто складывать ответы}

Если мы хотим измерить эмпатию, можно взять десять вопросов про эмпатию, просуммировать ответы и назвать это эмпатией. Это --- композитный балл; он полезен, но у него есть структурное ограничение: он предполагает, что все десять вопросов ловят одну и ту же вещь в одной и той же пропорции, и что шум измерения в них независим. В реальности часть дисперсии ответов --- общий конструкт (\emph{то, что все десять вопросов действительно ловят}), а часть --- специфичная для конкретной формулировки, конкретной ситуации, конкретного стиля ответа респондента. Латентный факторный анализ разделяет эти две доли явно. Фактор --- это \emph{ненаблюдаемая} переменная, восстановленная из паттерна корреляций между пунктами опросника: предполагается, что если эмпатия есть, то ответы на эмпатические вопросы должны коррелировать между собой сильнее, чем с ответами на вопросы про, скажем, психопатию. Сила связи каждого пункта с фактором называется \emph{нагрузкой} (factor loading); по сути, это «сколько процента этого пункта --- про эмпатию, а сколько --- про что-то ещё». Когда мы далее говорим о \emph{факторных оценках} (factor scores), мы имеем в виду оценку положения каждого респондента на этой ненаблюдаемой оси, полученную как взвешенная комбинация его ответов с весами, учитывающими величину нагрузок. В отличие от композита, факторная оценка автоматически \emph{пенализирует} шумные пункты и не раздувает корреляции между конструктами за счёт общей дисперсии (см. §\ref{sec:limitations} об артефакте мульти-шкальных весов в композитном счёте).

\subsubsection{Подтверждающий факторный анализ; косоугольная и ортогональная спецификации}

В подтверждающем факторном анализе (CFA) мы \emph{заранее} постулируем, какие пункты грузят на какой фактор, и проверяем, держится ли эта структура на данных. В противоположность эксплоративному анализу (EFA), где факторы находятся «снизу вверх» из корреляционной матрицы, CFA работает «сверху вниз»: у нас есть теория, у нас есть данные, и мы смотрим, совместимы ли они. Результат --- набор индексов соответствия, о которых ниже.

\emph{Косоугольная vs ортогональная спецификация} --- это решение о том, разрешаем ли мы факторам коррелировать между собой. Ортогональная CFA фиксирует все межфакторные корреляции в ноль --- удобно математически, но в психологии почти всегда неверно: эмпатия и психопатия \emph{реально} коррелируют отрицательно, ID и CA \emph{реально} коррелируют положительно, и заставлять эти связи быть нулевыми --- значит заранее искажать данные ради простоты модели. Косоугольная CFA, которую мы используем, разрешает корреляции между факторами и оценивает их непосредственно. Это содержательно: межфакторные корреляции --- часть результата, а не помеха.

\subsubsection{Индексы соответствия: что они говорят и чего не говорят}

Когда мы пишем «CFI = 0.950, RMSEA = 0.068, $\chi^2/df = 4.30$», за каждым числом стоит конкретная интуиция.

\textbf{$\chi^2$-статистика} измеряет, насколько наблюдаемая ковариационная матрица пунктов опросника расходится с той, которую предсказывает модель. Если модель идеальна, расхождение нулевое; чем больше расхождение, тем больше $\chi^2$. Беда в том, что $\chi^2$ растёт с размером выборки: при $N = 1426$ даже очень хорошая модель почти всегда отвергается по строгому $\chi^2$-тесту. Поэтому в современной психометрии $\chi^2$ читают как описание, а не как тест, и дополняют его относительными индексами.

\textbf{CFI (индекс сравнительного соответствия)} --- сравнение модели с «нулевой»: моделью, в которой все пункты независимы друг от друга (ни одного фактора вообще). CFI = 0.950 означает, что наша модель закрывает 95\% расстояния между нулевой моделью и идеальной. Принятые пороги: $\geq 0.95$ хорошо, $\geq 0.90$ приемлемо. CFI гораздо менее чувствителен к размеру выборки, чем $\chi^2$.

\textbf{RMSEA (средняя ошибка аппроксимации)} --- средняя ошибка модели в расчёте на одну степень свободы. Интуиция: насколько модель \emph{в среднем на параметр} промахивается мимо данных. Принятые пороги: $\leq 0.05$ хорошо, $\leq 0.08$ приемлемо. RMSEA = 0.068 --- приемлемо для сложной модели с реальными данными.

\textbf{$\chi^2/df$} --- отношение $\chi^2$-статистики к числу степеней свободы; грубая проверка «настолько ли плохо, на сколько ожидалось бы по случаю». $\leq 5$ приемлемо для больших выборок.

Важное предупреждение, которое мы повторяем в тексте: эти индексы работают в предположении многомерной нормальности данных; наши данные её нарушают (Mardia $p < 0.0001$), поэтому мы приводим и стандартные индексы при оценивании методом максимального правдоподобия (ML), и робастные, скорректированные Саторрой--Бентлером, а содержательные выводы всегда опираем на латентные факторные оценки, которые к этим коррекциям менее чувствительны.

\subsubsection{Бифакторная модель: общий фактор плюс остаточные специфики}

В §\ref{subsec:id_bifactor} мы применяем \emph{бифакторную} модель к пяти wave-1 шкалам пассивного блокирования: \textbf{SUP} (Suppression, подавление), \textbf{DIS} (Dissociation, диссоциация), \textbf{AIS} (Affective Isolation Scale, аффективная изоляция), \textbf{FLA} (Flattening, аффективное уплощение) и \textbf{IDS} (Internal Dissonance Scale, внутренний диссонанс --- историческое «родительское» название, от которого итоговый латент получил имя ID). Это особая архитектура CFA, в которой один общий фактор (ID\textsubscript{g}) грузит \emph{все} пункты опросника, а дополнительно к нему --- ортогональные общему фактору \emph{специфические} факторы (по одному на каждую подшкалу), грузящие только «свои» пункты. Это позволяет алгебраически развести дисперсию: сколько её поглощается общим конструктом, а сколько остаётся в специфических для подшкалы остатках. Метрика \textbf{ECV (доля объяснённой общей дисперсии)} $= 0.571$ говорит, что 57\% общей факторной дисперсии объясняется общим ID; остальные 43\% сидят в специфических подфакторах. ECV для отдельной подшкалы --- та же метрика, вычисленная внутри одной подшкалы: у SUP (подавления) она 0.31, то есть подавление в значительной степени \emph{не} про общий ID, а про что-то своё. Значение этого анализа для волны~2 --- структурное основание для редизайна пунктов.

\subsubsection{Медиация, модерация, параллельная медиация}

\emph{Медиация} отвечает на вопрос «через что». Если CA предсказывает \EAF{}, но мы предполагаем, что это предсказание идёт не напрямую, а через ID, --- мы тестируем медиацию. Формально: путь $a$ (CA $\to$ ID), путь $b$ (ID $\to$ \EAF{} при контроле CA), прямой путь $c'$ (CA $\to$ \EAF{} при контроле ID). Непрямой эффект --- произведение $a \cdot b$; суммарный эффект $= a \cdot b + c'$. Статистическая значимость непрямого эффекта тестируется через 95\%-ный бутстрэповский доверительный интервал: если интервал не включает ноль, медиация значима.

\emph{Модерация} отвечает на другой вопрос --- «при каких условиях». Модератор --- переменная, которая меняет \emph{силу} связи между двумя другими переменными. В регрессии это член взаимодействия: ID $\times$ A в уравнении \EAF{}. Если коэффициент при взаимодействии значим, сила связи ID $\to$ \EAF{} меняется в зависимости от уровня A.

\emph{Параллельная медиация} --- расширение: один предиктор (CA), один исход (\EAF{}), несколько медиаторов (ID, P, M), оценённых \emph{одновременно}. Каждый получает свой непрямой эффект (через себя), и мы можем тестировать, какие из них значимы, а какие пусты. В нашем случае канал CA$\to$P$\to$\EAF{} оказался пуст: он существует формально, но не несёт дисперсии. Это содержательно: P архитектурно ортогонален детской истории.

\subsubsection{Двойно-модерированная медиация и индекс \IMM}

Двойно-модерированная медиация --- когда \emph{оба} участка канала медиации ($a$ и $b$) модерируются одной и той же переменной. В нашей центральной спецификации A модерирует и CA$\to$ID (первую стадию), и ID$\to$\EAF{} (вторую стадию). Вопрос: объединяются ли эти модерации в согласованную картину или они независимы?

\textbf{Индекс двойно-модерированной медиации (\IMM)} \citep{hayes2018} --- произведение двух коэффициентов взаимодействия: $a_{\text{mod}} \times b_{\text{mod}}$. Интуитивно это «на сколько меняется непрямой эффект при единичном изменении модератора на обеих стадиях сразу». Если \IMM{} значим, канал как целое перестраивается под влиянием модератора, а не просто получает дополнительный шум. Тестируется через бутстрэповский доверительный интервал.

Ключевая интуиция для читателя, привыкшего к простому взаимодействию: обычная модерация отвечает на вопрос «зависит ли сила связи от условия»; \IMM{} отвечает на вопрос «зависит ли \emph{вся цепь} от условия одинаковым, согласованным образом». Это сильнее обычной модерации и слабее утверждения о полностью перестроенной архитектуре. Рис.~\ref{fig:imm_intuition} показывает это графически.

\subsubsection{Рефлексивные и формативные конструкты: два способа собирать содержательную единицу}

Это различение, на котором держится вся дискуссия о A (§\ref{subsec:A_as_field_condition}, §\ref{sec:limitations}). В классической психометрической модели конструкт считается \emph{рефлексивным}: он --- причина своих индикаторов. Эмпатия «вызывает» высокие ответы на эмпатические пункты; если у вас действительно есть эмпатия, это проявится в нескольких связанных поведениях. Из этого следуют три предсказания: индикаторы должны положительно коррелировать между собой, должны иметь положительные нагрузки на общий фактор, и удаление одного индикатора не меняет смысла конструкта.

\emph{Формативный} конструкт устроен обратно: индикаторы не следствия, а \emph{составные части} конструкта. Классический пример --- социально-экономический статус: доход, образование, профессия не симптомы СЭС, а \emph{компоненты}, из которых СЭС складывается. Убрать доход из модели СЭС --- значит изменить сам конструкт. Индикаторы формативного конструкта не обязаны положительно коррелировать между собой, а нагрузки могут быть \emph{биполярными}: если одна компонента положительна, а другая отрицательна, это не противоречие, а структурное свойство.

В §\ref{sec:results} и §\ref{sec:limitations} мы показываем, что A в текущем пуле пунктов ведёт себя как формативный композит, а не как рефлексивный латент: CFA даёт биполярные нагрузки, и подстановка \emph{латентного} A в уравнение условной медиации переворачивает знаки ключевых коэффициентов. Это не провал анализа, а свидетельство того, что инструмент измеряет суперпозицию двух разных вещей (наблюдающее отношение и травматическая гипервнимательность), которые теоретически взвешенный композит разделяет, а голая CFA --- нет. Отсюда задача волны~2: развести эти два конструкта на уровне пунктов опросника.

\subsubsection{Бутстрэповский интервал и байесовский постериор: два способа выразить неопределённость}

\textbf{Бутстрэповский доверительный интервал} --- непараметрический способ построения доверительного интервала для статистики, чьё выборочное распределение неизвестно или не нормально. Процедура: из исходной выборки $N$ раз с возвращением вынимаем бутстрэп-выборку того же размера, пересчитываем на ней статистику, записываем. После 2000--5000 итераций у нас есть эмпирическое распределение статистики; 2{,}5-я и 97{,}5-я перцентили этого распределения --- границы 95\%-ного интервала. Применяется к непрямым эффектам и \IMM, потому что их выборочное распределение --- произведение нормальных, то есть \emph{не} нормальное.

\textbf{Байесовский постериор} --- другая логика. Вместо доверительного интервала, в который «попадает истинное значение с вероятностью 95\%», байесовский подход даёт \emph{апостериорное распределение} параметра: после того, как мы увидели данные, насколько вероятно каждое значение параметра. Постериор $=$ приор (наши убеждения до данных) $\times$ правдоподобие (вероятность данных при каждом значении параметра), пронормированный. 95\%-ный интервал наивысшей плотности (HDI) --- самый узкий интервал, содержащий 95\% вероятностной массы. $P(\IMM < 0) = 0{,}97$ означает буквально: с учётом данных и приора вероятность того, что \IMM{} отрицателен, равна 97\%. Эта формулировка содержательно богаче, чем «интервал не включает ноль», и устойчива к разным формам приоров, что мы и демонстрируем в §\ref{subsubsec:bayesian_imm}.

\subsubsection{Что это всё вместе даёт}

Цепочка анализов, которую читатель увидит в §\ref{sec:results}, следует одной логике: \textbf{(1)}~сначала мы устанавливаем, что шесть факторов различимы (шестифакторная косоугольная CFA с приемлемыми индексами соответствия); \textbf{(2)}~затем --- что в некоторых парах они ортогональны (частная корреляция, ID$\perp$P); \textbf{(3)}~затем тестируем причинный каскад, который эти факторы формируют (параллельная медиация и условная медиация с двойной модерацией); \textbf{(4)}~проверяем устойчивость этого каскада (байесовский пересчёт); \textbf{(5)}~собираем всё в одно объединённое структурное уравнение, где видно, как каналы реорганизуются при одновременной оценке; \textbf{(6)}~различаем топологические подгруппы (первичная и вторичная психопатия, интегрированный выживший); \textbf{(7)}~проверяем, действительно ли ключевой модератор (A) работает как рефлексивный латент, и обнаруживаем, что нет, --- отсюда приоритет волны~2. Каждый следующий шаг --- не независимый тест, а уточнение предыдущей картины.

\subsection{Статистический подход}
\label{subsec:stat_approach}

\subsubsection{Инструменты}

Анализ выполнен в Python 3.10. Ключевые библиотеки: \texttt{pandas}, \texttt{numpy}, \texttt{scipy}, \texttt{scikit-learn}, \texttt{semopy} (CFA, SEM), \texttt{factor-analyzer} (EFA), \texttt{pymc} (байесовская модель для IMM). Фактор-скоры и conditional process SEM вычислены в \texttt{semopy}. Bootstrap 95\% CI (2000--5000 draws в зависимости от модели) использован для всех indirect effects и IMM.

Fit indices отчитываются в двух формах: (а)~классический ML (нормально-теоретический) и (б)~Satorra--Bentler T$_1$ масштабированная $\chi^2$-статистика через \texttt{semopy.stats.calc\_chi2\_sb} с построением четвёрто-моментной весовой матрицы (\texttt{prepare\_wls}) поверх ML-посадки. Для 6-факторной oblique CFA на split-B ($n = 713$): ML $\chi^2(120) = 516.04$, CFI $= 0.950$, TLI $= 0.936$, RMSEA $= 0.068$; SB-масштабированные $\chi^2_{\mathrm{SB}}(120) = 934.25$, CFI$_{\mathrm{SB}} = 0.873$, TLI$_{\mathrm{SB}} = 0.838$, RMSEA$_{\mathrm{SB}} = 0.098$, scaling factor $c = 0.552$. Нетипичное $c < 1$ (ML-статистика дефлирована относительно SB-корректированной) и существенный спад CFI при переходе к SB указывают на то, что reviewer-стандарт требует параллельной реализации в \texttt{lavaan} R (MLR, WLSMV на сырых Likert-айтемах); это отмечено в §\ref{sec:limitations} как инфраструктурный gap. Содержательные выводы опираются на latent factor scores, извлечённые из ML-посадки, а не на $\chi^2$ per se, и в этой части устойчивы к robust/classical различию.

\subsubsection{Conditional process SEM с A как модератором обеих стадий}
\label{subsubsec:cp_sem}

Центральный анализ --- moderated-moderated mediation канала CA~$\to$~ID~$\to$~\EAF{} с A как модератором обеих стадий. Формальная модель:
\begin{align}
\text{ID} &= \beta_0 + a \cdot \text{CA} + \beta_{A_{\text{ID}}} \cdot A + a_{\text{mod}} \cdot (\text{CA} \times A) + \varepsilon_1 \\
\text{\EAF{}} &= \beta_0 + b \cdot \text{ID} + c' \cdot \text{CA} + \beta_{A_{\text{E}}} \cdot A + b_{\text{mod}} \cdot (\text{ID} \times A) + \varepsilon_2
\end{align}
Индекс моделируемой модерации: $\IMM = a_{\text{mod}} \times b_{\text{mod}}$. \IMM{} тестируется через bootstrap 95\% CI; ненулевое CI --- свидетельство значимой модерации обеих стадий одновременно.

\subsubsection{\texorpdfstring{Параллельная медиация CA $\to$ \{ID, P, M\} $\to$ \EAF{}}{Параллельная медиация CA -> {ID, P, M} -> E\_AF}}

Прямой тест двухканальной архитектуры: единый SEM с CA как предиктором, \EAF{} как исходом, и ID, P, M как тремя параллельными медиаторами. Estimated на latent factor scores split-B ($n = 710$). Bootstrap 5000 итераций. Отчитываются: a-пути (CA $\to$ медиатор), b-пути (медиатор $\to$ \EAF{} при контроле остальных), indirect effects через каждый медиатор с CI, contrast tests между парами indirect effects, direct c'.

\subsubsection{Unified SEM: все пути одновременно}

Наиболее полная спецификация объединяет параллельную медиацию (§\ref{subsubsec:cp_sem}) с A-модерацией. CA предсказывает ID (с A как главным эффектом и CA$\times$A как модератором), P и M (оба линейно от CA без A-модерации). \EAF{} регрессируется на все три медиатора + A (главный эффект) + ID$\times$A (b-модерация) + c' (прямой CA-эффект). 2000 bootstrap draws. Nested comparison (M0: E~$\sim$~CA; M1: +\{ID, P, M\}; M2: +A main; M3: +ID$\times$A) через $F$-, AIC- и BIC-тесты.

\subsubsection{Bayesian reanalysis ключевого IMM}

В предвосхищение reviewer-скептицизма относительно frequentist CI, касающегося нуля в верхней границе, мы повторили conditional process модель (§\ref{subsubsec:cp_sem}) в байесовской спецификации через PyMC, NUTS sampler, 4 цепи $\times$ (2000 tune + 4000 draws) = 16000 post-warmup draws, все $\hat{r} = 1.00$. Три приора сравниваются: flat (Normal(0, 100) на всех коэффициентах), weakly informative (Normal(0, 1)), sceptical (Normal(0, 0.1) на обеих интеракциях $a_3$, $b_5$). Отчитываются posterior mean, 95\% HDI, $P(\IMM < 0)$.

\subsubsection{Two-tier A-CFA (Tier 5 и 5b)}

Подтверждение формативной природы A выполнено через поэтапную A-CFA.
\begin{description}
    \item[Tier 5] --- три варианта single-factor reflective CFA: (a) full model (все 31 A-взвешенный айтем); (b) core model (айтемы с $|w_A| \geq 0.3$, $n = 18$); (c) purest model (17 метакогнитивных self-observation айтемов, отобранных по содержанию).
    \item[Tier 5b] --- two-factor CFA на purest-17 после обнаружения биполярных нагрузок: A\_presence (8 items: позитивная феноменология) vs A\_absence (9 items: феноменология диссоциации), correlated factors.
\end{description}
Обе A-CFA процедуры запускаются параллельно с подстановкой полученных factor scores в conditional process SEM для проверки, воспроизводит ли latent-A сигнал, полученный на композите.

\subsubsection{Валидация supplementary}

\begin{itemize}
    \item \textbf{Split-half:} EFA на A-половине ($N = 713$), CFA на B-половине ($N = 713$), seed = 42 для воспроизводимости.
    \item \textbf{Gender invariance:} configural $\to$ metric $\to$ scalar через chain-constrained CFA с sequential model comparison ($\Delta$CFI, $\Delta$RMSEA criteria).
    \item \textbf{Partial correlations} с Lie-шкалой как ковариатой --- контроль socially desirable responding.
    \item \textbf{Network psychometrics} (graphical lasso): partial correlation network на шести факторах, $\gamma = 0.5$ regularization.
    \item \textbf{ROC analyses:} predictive performance каждого фактора для self-reported clinical diagnoses, с 95\% AUC CI через DeLong method.
\end{itemize}

\subsection{Деривация дополнительных конструктов}

\subsubsection{\texorpdfstring{Primary vs secondary psychopathy (topological $2\times 2$)}{Primary vs secondary psychopathy (topological 2x2)}}

Rule-based классификация, мотивированная \citet{karpman1941}: участник считается \emph{primary-like} при P > медианы, \EAF{} < медианы, ID < медианы, CA < медианы; \emph{secondary-like} при P > медианы, \EAF{} < медианы, ID > медианы, CA > медианы. Без prior claim на categorical kinds --- классификация технически dimensional-based thresholding.

\subsubsection{\texorpdfstring{Integrated survivor (hi-CA $\times$ hi-A)}{Integrated survivor (hi-CA x hi-A)}}

Участник классифицируется как integrated-survivor-like при CA > 75-й перцентили и A > 75-й перцентили. Этот порог строже median-split для выделения выраженного профиля; результаты с median-split доступны в supplementary.

\subsection{Transparency и pre-registration}

Текущий анализ post-hoc относительно первоначального дизайна инструмента. Мы эксплицитно это маркируем в каждом результате, где post-hoc нюанс релевантен. Pre-registered replication на новой выборке --- обязательный следующий шаг (описан в \texttt{PROGRAM\_STRUCTURE.md}). Все скрипты анализа и intermediate CSV-файлы выложены в \texttt{analysis\_output/fresh\_analysis/}; анонимизированный raw dataset --- в \texttt{data/anonymized/}.
